\documentclass[a4paper,10pt]{exam}
\usepackage[T1]{fontenc}
\usepackage[utf8]{inputenc}
\usepackage[italian]{babel}

% --- TIPOGRAFIA: stessa del modello canonico ---
\usepackage{newpxtext}
\usepackage{newpxmath}
\usepackage{microtype}

\usepackage{amsfonts}
\usepackage{amsmath}
\usepackage{graphicx}
\usepackage{siunitx}
\usepackage{physics}
\usepackage{geometry}
\usepackage{multicol}
\usepackage[table]{xcolor}
\usepackage{tcolorbox}
\usepackage{tikz}
\usepackage{qrcode}

\AtBeginDocument{\RenewCommandCopy\qty\SI}
\newcommand{\um}[2]{\SI[output-decimal-marker={,}]{#1}{#2}}

\geometry{top=1.6cm,bottom=1.0cm,left=0.8cm,right=0.8cm,headsep=0.3cm}

% INTESTAZIONE: solo punteggi e codice variante (deroga al divieto di fila, solo per i quiz).
\pagestyle{headandfoot}
\firstpageheader{}{\small Punti per risposta corretta: 4 $\cdot$ non data: 1 $\cdot$ errata: 0}{\scriptsize Codice: \textbf{5ZAN}}
\runningheader{}{\small Punti per risposta corretta: 4 $\cdot$ non data: 1 $\cdot$ errata: 0}{\scriptsize Codice: \textbf{5ZAN}}
\firstpagefooter{}{\scriptsize Pagina \thepage\ di \numpages}{}
\runningfooter{}{\scriptsize Pagina \thepage\ di \numpages}{}

% --- STILE MONOCROMATICO ---
\newcommand{\badgeQ}[1]{%
  \tikz[baseline=(char.base)]{%
    \node[shape=rectangle, rounded corners=2pt, fill=gray!15, draw=gray!40, line width=0.4pt, inner sep=2.5pt, minimum width=1.6em] (char) {\small\textbf{#1}};%
  }%
}
\renewcommand{\questionlabel}{\badgeQ{\thequestion}\enspace}
\renewcommand{\choicelabel}{\textbf{\Alph{choice}}\enspace}
\newcommand{\cbox}{\ensuremath{\vcenter{\hbox{\tikz{\draw[black!75, line width=0.5pt, rounded corners=1.2pt] (0,0) rectangle (3.0mm, 3.0mm);}}}}}

\tcbset{
  istruzionibox/.style={
    colback=gray!6, colframe=gray!30, arc=2pt, boxrule=0.5pt,
    left=8pt, right=8pt, top=5pt, bottom=5pt, fontupper=\small
  }
}

\begin{document}
\begin{tcolorbox}[istruzionibox]
È consentito l'uso di penne nere o blu (non è consentito l'uso del lapis e del bianchetto).
È consentito l'uso della calcolatrice scientifica non grafica.
Per ogni domanda una sola risposta è corretta. Le risposte si danno sul foglio risposte (ultima pagina): annerire completamente il cerchio scelto; per annullare una risposta annerire due o più cerchi. Le correzioni non sono ammesse e valgono come risposta non data.
\end{tcolorbox}

\vspace{2mm}
\noindent
\textbf{Cognome e Nome:}\enspace\hrulefill\hspace{1cm}%
\textbf{Classe:}\enspace\makebox[2.5cm]{\hrulefill}\hspace{1cm}%
\textbf{Data:}\enspace\makebox[2.5cm]{\hrulefill}

\vspace{2mm}
\begin{questions}
\question Un carrello di $\um{3}{kg}$ a $\um{4}{m/s}$ si aggancia a uno fermo di $\um{1}{kg}$. La velocità finale è:
\vspace{0.2em}
\begin{multicols}{2}
\begin{choices}
\choice $\um{1}{m/s}$
\choice $\um{4}{m/s}$
\choice $\um{2}{m/s}$
\choice $\um{3}{m/s}$
\end{choices}
\end{multicols}

\question Nel pendolo balistico l'energia meccanica si conserva:
\vspace{0.2em}
\begin{choices}
\choice mai
\choice durante l'urto
\choice prima dell'urto soltanto
\choice solo dopo l'urto
\end{choices}

\question L'energia cinetica dissipata nell'urto completamente anelastico precedente è sempre:
\vspace{0.2em}
\begin{choices}
\choice uguale all'energia finale
\choice sempre nulla
\choice positiva o nulla
\choice negativa
\end{choices}

\question Un corpo $m$ urta elasticamente un corpo fermo di massa uguale. Dopo l'urto:
\vspace{0.2em}
\begin{choices}
\choice entrambi si fermano
\choice il primo si ferma e il secondo prende la sua velocità
\choice il primo rimbalza con la stessa velocità
\choice i due procedono uniti
\end{choices}

\question Un proiettile di $\um{10}{g}$ a $\um{300}{m/s}$ si conficca in un blocco fermo di $\um{2.99}{kg}$. La velocità subito dopo è:
\vspace{0.2em}
\begin{multicols}{2}
\begin{choices}
\choice $\um{0.1}{m/s}$
\choice $\um{1}{m/s}$
\choice $\um{300}{m/s}$
\choice $\um{10}{m/s}$
\end{choices}
\end{multicols}

\question In un urto elastico si conserva, oltre alla quantità di moto:
\vspace{0.2em}
\begin{choices}
\choice la velocità del centro di massa e basta
\choice l'energia cinetica totale
\choice nulla
\choice la massa di ciascun corpo
\end{choices}

\question Dopo un urto elastico in una dimensione fra masse diverse, la velocità relativa dei due corpi:
\vspace{0.2em}
\begin{choices}
\choice cambia segno e mantiene il modulo
\choice si dimezza
\choice resta uguale
\choice si annulla
\end{choices}

\question Quale grandezza ha le stesse dimensioni della quantità di moto?
\vspace{0.2em}
\begin{choices}
\choice la forza
\choice la potenza
\choice l'impulso
\choice l'energia
\end{choices}

\end{questions}

\clearpage\ifodd\value{page}\thispagestyle{empty}\null\clearpage\fi
\thispagestyle{empty}\null
\begin{tikzpicture}[remember picture,overlay,shift={(current page.north west)},x=1mm,y=-1mm]
\fill (7.0,7.0) rectangle (17.0,17.0);
\fill (195.0,9.0) rectangle (201.0,15.0);
\fill (195.0,282.0) rectangle (201.0,288.0);
\fill (9.0,282.0) rectangle (15.0,288.0);
\node[anchor=west] at (24,20) {\large\textbf{Foglio risposte}};
\node[anchor=west] at (24,25) {\small Codice: \textbf{5ZAN}};
\node[anchor=west] at (20,31) {\scriptsize Cognome e Nome (stampatello)};
\draw[line width=0.5pt] (20,33) rectangle (138,45);
\node[anchor=north west,inner sep=0] at (150,18) {\qrcode[height=36mm,level=M]{QZ-5ZAN}};
\node[anchor=north west,text width=125mm,align=justify] at (20,49) {\scriptsize Annerisci completamente il cerchio della risposta scelta con penna nera o blu. Per annullare una risposta annerisci due o più cerchi della domanda. Le correzioni non sono ammesse: cerchi cancellati, coperti o anneriti in parte possono rendere la domanda non data. Non fare segni fuori dai cerchi.};
\node at (30,77) {\small\textbf{A}};
\node at (38,77) {\small\textbf{B}};
\node at (46,77) {\small\textbf{C}};
\node at (54,77) {\small\textbf{D}};
\node at (82,77) {\small\textbf{A}};
\node at (90,77) {\small\textbf{B}};
\node at (98,77) {\small\textbf{C}};
\node at (106,77) {\small\textbf{D}};
\node at (134,77) {\small\textbf{A}};
\node at (142,77) {\small\textbf{B}};
\node at (150,77) {\small\textbf{C}};
\node at (158,77) {\small\textbf{D}};
\fill[black!7] (14,80.0) rectangle (59,88.0);
\node[anchor=east] at (24,84) {\small\textbf{1}};
\draw[line width=0.5pt] (30,84) circle (2.3mm);
\draw[line width=0.5pt] (38,84) circle (2.3mm);
\draw[line width=0.5pt] (46,84) circle (2.3mm);
\draw[line width=0.5pt] (54,84) circle (2.3mm);
\node[anchor=east] at (24,92) {\small\textbf{2}};
\draw[line width=0.5pt] (30,92) circle (2.3mm);
\draw[line width=0.5pt] (38,92) circle (2.3mm);
\draw[line width=0.5pt] (46,92) circle (2.3mm);
\draw[line width=0.5pt] (54,92) circle (2.3mm);
\fill[black!7] (14,96.0) rectangle (59,104.0);
\node[anchor=east] at (24,100) {\small\textbf{3}};
\draw[line width=0.5pt] (30,100) circle (2.3mm);
\draw[line width=0.5pt] (38,100) circle (2.3mm);
\draw[line width=0.5pt] (46,100) circle (2.3mm);
\draw[line width=0.5pt] (54,100) circle (2.3mm);
\fill[black!7] (66,80.0) rectangle (111,88.0);
\node[anchor=east] at (76,84) {\small\textbf{4}};
\draw[line width=0.5pt] (82,84) circle (2.3mm);
\draw[line width=0.5pt] (90,84) circle (2.3mm);
\draw[line width=0.5pt] (98,84) circle (2.3mm);
\draw[line width=0.5pt] (106,84) circle (2.3mm);
\node[anchor=east] at (76,92) {\small\textbf{5}};
\draw[line width=0.5pt] (82,92) circle (2.3mm);
\draw[line width=0.5pt] (90,92) circle (2.3mm);
\draw[line width=0.5pt] (98,92) circle (2.3mm);
\draw[line width=0.5pt] (106,92) circle (2.3mm);
\fill[black!7] (66,96.0) rectangle (111,104.0);
\node[anchor=east] at (76,100) {\small\textbf{6}};
\draw[line width=0.5pt] (82,100) circle (2.3mm);
\draw[line width=0.5pt] (90,100) circle (2.3mm);
\draw[line width=0.5pt] (98,100) circle (2.3mm);
\draw[line width=0.5pt] (106,100) circle (2.3mm);
\fill[black!7] (118,80.0) rectangle (163,88.0);
\node[anchor=east] at (128,84) {\small\textbf{7}};
\draw[line width=0.5pt] (134,84) circle (2.3mm);
\draw[line width=0.5pt] (142,84) circle (2.3mm);
\draw[line width=0.5pt] (150,84) circle (2.3mm);
\draw[line width=0.5pt] (158,84) circle (2.3mm);
\node[anchor=east] at (128,92) {\small\textbf{8}};
\draw[line width=0.5pt] (134,92) circle (2.3mm);
\draw[line width=0.5pt] (142,92) circle (2.3mm);
\draw[line width=0.5pt] (150,92) circle (2.3mm);
\draw[line width=0.5pt] (158,92) circle (2.3mm);
\end{tikzpicture}
\end{document}
